\section{Divergences and Dual Norms}
\label{sec-divergences-dual-norms}
\subsection{Dual norms (Integral Probability Metrics)}
\label{sec-dual-norms}
\paragraph{Integral probability metrics.}
\begin{defn}[Dual seminorm and integral probability metric]\label{def-dual-norm-ipm}
Let $B$ be a nonempty symmetric convex class of measurable functions that are integrable against the measures under consideration. For a signed measure $\xi$, define the extended dual seminorm
\eql{\label{eq-dual-norm-cont}
\norm{\xi}_B \eqdef
\usup{\f} \enscond{ \int_{\Xx} \f(x) \d\xi(x) }{ \f \in B}.
}
It is a norm when it is finite and $B$ separates signed measures. When applied to $\al-\be$ for probability measures, it is called an integral probability metric (IPM); it is a genuine metric precisely when $B$ separates probability measures.
\end{defn}
\begin{prop}[Metrization by dual norms]\label{prop-dual-norm-metrization}
Assume that $\Xx$ is compact, that $B=-B$, and that the measures considered are probability measures.
\begin{enumerate}
\item If every function in $\Cc(\Xx)$ can be uniformly approximated by elements of $\Span(B)$, then $\norm{\al_n-\al}_B \rightarrow 0$ implies $\al_n \rightharpoonup \al$.
\item If $B \subset \Cc(\Xx)$ is compact for $\norm{\cdot}_\infty$, then $\al_n \rightharpoonup \al$ implies $\norm{\al_n-\al}_B \rightarrow 0$.
\end{enumerate}
\end{prop}
\begin{proof}
For the first implication, $\norm{\al_n-\al}_B\to0$ and the symmetry of $B$ imply
\(\abs{\dotp{f}{\al_n-\al}}\leq \norm{\al_n-\al}_B \qforq f\in B.\)
If $h=\sum_{j=1}^J c_j f_j\in\Span(B)$, then
\(\abs{\dotp{h}{\al_n-\al}} \leq \pa{\sum_{j=1}^J|c_j|}\norm{\al_n-\al}_B,\)
so integrals converge for every $h\in\Span(B)$. Let $u\in\Cc(\Xx)$ and choose such an $h$ with $\norm{u-h}_\infty\leq\eta$. Since $\al_n$ and $\al$ are probabilities,
\(\abs{\dotp{u}{\al_n-\al}} \leq \abs{\dotp{h}{\al_n-\al}}+2\eta.\)
Taking the limsup as $n\to\infty$ and then letting $\eta\to0$ gives $\dotp{u}{\al_n}\to\dotp{u}{\al}$ for all $u\in\Cc(\Xx)$, which is weak convergence.

For the second implication, assume $\al_n\rightharpoonup\al$ and choose a subsequence $(\al_{n_k})_k$ such that $\norm{\al_{n_k}-\al}_B\to\limsup_n\norm{\al_n-\al}_B$. Since $B$ is compact and the map $f\mapsto\dotp{f}{\al_{n_k}-\al}$ is continuous on $B$, the supremum is attained by some $f_{n_k}\in B$. Extracting a further subsequence if needed, $f_{n_k}\to f$ uniformly for some $f\in B$. Then
\(\dotp{f_{n_k}}{\al_{n_k}-\al} = \dotp{f}{\al_{n_k}-\al} +\dotp{f_{n_k}-f}{\al_{n_k}} -\dotp{f_{n_k}-f}{\al}.\)
The first term tends to zero by weak convergence and the last two by uniform convergence. Hence every limsup subsequence has limit zero, proving $\norm{\al_n-\al}_B\to0$.
\end{proof}
\begin{cor}[Wasserstein metrizes weak convergence]\label{cor-topol-wass}
On a compact metric space, $\Wass_p$ metrizes weak convergence on probability measures for every $p\geq1$.
\end{cor}
\begin{proof}
For $p=1$, take $B=\enscond{f}{\Lip(f)\leq1}$. The span of $B$ contains all Lipschitz functions, and Lipschitz functions are dense in $\Cc(\Xx)$ on compact metric spaces.

Conversely, constants do not change the pairing with $\al_n-\al$. Fix $x_0\in\Xx$ and normalize potentials by $f(x_0)=0$. The normalized unit Lipschitz ball is uniformly bounded by $\mathrm{diam}(\Xx)$ and equicontinuous, hence compact in $\norm{\cdot}_\infty$ by Arzel\`a--Ascoli. Proposition~\ref{prop-dual-norm-metrization} gives $\Wass_1(\al_n,\al)\to0$. Proposition~\ref{prop-comp-wass-p} shows that all $\Wass_p$ distances induce the same topology on a compact space, so the result follows for every $p\geq1$.
\end{proof}
\subsection{Dual RKHS Norms and Maximum Mean Discrepancies}
\label{sec-rkhs-mmd}
\begin{defn}[Positive and conditionally positive kernels]\label{def-positive-kernels}
A symmetric function $\Krkhs:\Xx\times\Xx\to\RR$ is positive definite if for every $n\geq1$, every $x_1,\ldots,x_n\in\Xx$ and every $r\in\RR^n$,
\eql{\label{eq-dual-kern}
\sum_{i,j=1}^n r_i r_j \Krkhs(x_i,x_j)\geq0.
}
It is conditionally positive definite if the same inequality is required only for zero-sum vectors, $\dotp{r}{\ones_n}=0$.
\end{defn}
\begin{defn}[Kernel seminorm and MMD]\label{def-kernel-mmd-norm}
Let $\Krkhs$ be positive definite. More generally, let $\Krkhs$ be conditionally positive definite and restrict attention to signed measures of total mass zero. For a signed measure $\xi$ with finite kernel energy, the associated seminorm is
\eql{\label{eq-kernel-dual}
\norm{\xi}^2_{\Krkhs}
\eqdef
\iint_{\Xx\times\Xx}\Krkhs(x,y)\d\xi(x)\d\xi(y).
}
For two probability measures, the maximum mean discrepancy associated with $\Krkhs$ is \(\MMD_{\Krkhs}(\al,\be)\eqdef\norm{\al-\be}_{\Krkhs}.\)
It is a genuine distance when the kernel is characteristic, meaning that $\norm{\al-\be}_{\Krkhs}=0$ implies $\al=\be$.
\end{defn}
\(\widetilde{\Krkhs}(x,y) =\Krkhs(x,y)-\Krkhs(x,x_0)-\Krkhs(x_0,y)+\Krkhs(x_0,x_0)\)
\begin{prop}[Kernel seminorm as an RKHS dual norm]\label{prop-kernel-rkhs-dual}
Let $\RKHS$ be the RKHS with reproducing kernel $\Krkhs$, and assume that the kernel mean embedding
\(m_\xi \eqdef \int \Krkhs(x,\cdot)\d\xi(x)\)
is well-defined for the signed measure $\xi$. Then
\(\norm{\xi}_{\Krkhs} = \sup_{\norm{h}_{\RKHS}\leq1} \int h(x)\d\xi(x),\)
so $\norm{\cdot}_{\Krkhs}$ is the dual seminorm in the sense of~\eqref{eq-dual-norm-cont} associated with the RKHS unit ball.
\end{prop}
\begin{proof}
By the reproducing property,
\[
\int h(x)\d\xi(x)
=
\left\langle h,\int \Krkhs(x,\cdot)\d\xi(x)\right\rangle_{\RKHS}
=
\langle h,m_\xi\rangle_{\RKHS}.
\]
Cauchy--Schwarz gives \(\sup_{\norm{h}_{\RKHS}\leq1}\int h\d\xi = \norm{m_\xi}_{\RKHS}.\)
Finally, \(\norm{m_\xi}_{\RKHS}^2 = \iint \Krkhs(x,y)\d\xi(x)\d\xi(y),\)
which is exactly~\eqref{eq-kernel-dual}.
\end{proof}
\begin{prop}[Universal kernels metrize weak convergence]\label{prop-mmd-metrization}
Assume that $\Xx$ is compact and that the RKHS generated by the continuous kernel $\Krkhs$ is dense in $\Cc(\Xx)$ for the uniform norm. Then
\[
\MMD_{\Krkhs}(\al_n,\al)\to0
\quad\Longleftrightarrow\quad
\al_n\rightharpoonup\al
\]
for probability measures on $\Xx$.
\end{prop}
\begin{proof}
If $\MMD_{\Krkhs}(\al_n,\al)\to0$, then for every $g\in\RKHS$, \(\abs{\int g\d(\al_n-\al)} \leq \norm{g}_{\RKHS}\MMD_{\Krkhs}(\al_n,\al)\to0.\)
For any $h\in\Cc(\Xx)$ and any $\eta>0$, choose $g\in\RKHS$ with $\norm{h-g}_\infty\leq\eta$. Since $\al_n$ and $\al$ are probabilities,
\[
\left|\int h\,\d(\al_n-\al)\right|
\leq
2\eta+\left|\int g\,\d(\al_n-\al)\right|,
\]
and the last term tends to zero. This proves weak convergence. Conversely, if $\al_n\rightharpoonup\al$, then $\al_n\otimes\al_n$, $\al_n\otimes\al$ and $\al\otimes\al$ converge weakly on the compact product space. Applying this to the continuous bounded function $\Krkhs$ in the identity
\(\MMD_{\Krkhs}(\al_n,\al)^2 = \iint \Krkhs\,\d\al_n\d\al_n -2\iint \Krkhs\,\d\al_n\d\al +\iint \Krkhs\,\d\al\d\al\)
gives convergence to zero.
\end{proof}
In the special case where $\al$ is a discrete measure, one thus has the simple expression

\eq{
\norm{\al}_{\Krkhs}^2 = \sum_{i=1}^n \sum_{i'=1}^n \a_i\a_{i'} \KrkhsD_{i,i'} = \dotp{\KrkhsD\a}{\a}
\qwhereq
\KrkhsD_{i,i'} \eqdef \Krkhs(x_i,x_{i'}).
}
\eql{\label{eq-mmd-discr}
\norm{\al-\be}_{\Krkhs}^2 =
\sum_{i,i'} \a_i \a_{i'} \Krkhs(x_i,x_{i'})	+
\sum_{j,j'} \b_j \b_{j'} \Krkhs(y_j,y_{j'}) - 2
\sum_{i,j} \a_i \b_j \Krkhs(x_i,y_j).
}
\subsection{\texorpdfstring{$\phi$-divergences}{phi-divergences}}
\label{sec-phi-div}
\paragraph{Definition by density ratios.}
\begin{defn}[Entropy function]
\label{def_entropy}
A function $\phi: \RR \to \RR \cup \{+\infty\}$ is an entropy function if it is proper, lower semicontinuous and convex, with $\dom \phi\subset [0,+\infty)$ and $\dom \phi \cap (0,+\infty) \neq \emptyset$. Its recession slope is
\eq{
\phi'_\infty = \lim_{s\rightarrow +\infty} \frac{\phi(s)}{s} \in \RR \cup \{+\infty\}.
}
\end{defn}
If $\phi'_\infty = \infty$, then $\phi$ grows faster than any linear function and $\phi$ is said to be \emph{superlinear}.

\begin{defn}[$\phi$-Divergences]
\label{def_divergence}
Let $\phi$ be an entropy function and let $\al,\be\in\Mm_+(\Xx)$. Write the Lebesgue decomposition of $\al$ with respect to $\be$ as
$\al=(\d\al/\d\be)\be+\al^\perp$, where $\al^\perp\perp\be$. The associated divergence is
\eql{\label{eq-phi-div}
\Divergm_\phi (\al|\be) \eqdef \int_\Xx \phi\left(\frac{\d \al}{\d \be} \right) \d \be
+ \phi'_\infty \al^{\perp}(\Xx),
}
with the convention $0\cdot(+\infty)=0$. It is extended by $+\infty$ whenever either argument is not a nonnegative measure.
\end{defn}
\eql{\label{eq-div-disc-meas}
\al=\sum_i \a_i \de_{x_i}
\qandq \be=\sum_i \b_i \de_{x_i}
}
\eql{\label{eq-discr-diverg}
\DivergmD_\phi(\a|\b) = \sum_{i \in \Supp(\b)} \phi\pa{ \frac{\a_i}{\b_i} } \b_i + \phi'_\infty \sum_{i \notin \Supp(\b)} \a_i,
}
where $\Supp(\b) \eqdef \enscond{i \in \range{n}}{ \b_i \neq 0 }$.

\begin{proposition}[Basic properties of $\phi$-divergences]\label{prop-basic-phi-divergence-properties}
If $\phi$ is an entropy function, then $\Divergm_\phi$ is jointly $1$-homogeneous, convex and weak-* lower semicontinuous in $(\al,\be)$.
\end{proposition}
\begin{proof}
Introduce the lower-semicontinuous perspective
\eq{
\foralls (u,v) \in (\RR_+)^2, \quad
\psi(u,v) = \choice{
v\phi(u/v) \qifq v>0 \\
u \phi'_\infty \qifq v=0.
}
}
The joint $1$-homogeneity follows directly from this formula. To see convexity, first take $v_0,v_1>0$ and $\tau\in[0,1]$. With $\bar v=(1-\tau)v_0+\tau v_1$ and
\(\theta_0=\frac{(1-\tau)v_0}{\bar v}, \qquad \theta_1=\frac{\tau v_1}{\bar v},\)
one has $\theta_0+\theta_1=1$ and
\(\frac{(1-\tau)u_0+\tau u_1}{\bar v} = \theta_0\frac{u_0}{v_0}+\theta_1\frac{u_1}{v_1}.\)
Jensen's inequality for $\phi$, followed by multiplication by $\bar v$, proves convexity of $\psi$ when both second arguments are positive. Its recession value at $v=0$ gives the lower-semicontinuous convex extension to the closed quadrant. In the discrete case,
\(\DivergmD_\phi(\a|\b) = \sum_{i} \psi(\a_i,\b_i),\)
so homogeneity and convexity follow by summation. For general finite measures, representing both arguments with respect to a common dominating measure yields the integral of $\psi$; the standard lower-semicontinuity theorem for convex integral functionals gives weak-* lower semicontinuity, with the recession value accounting for singular mass.
\end{proof}
\begin{proposition}[Non-negativity of $\phi$-divergences]\label{phi-div-positive}
Assume that $\phi$ is normalized by $\phi(1)=0$. For probability distributions $(\al,\be) \in \Mm_+^1(\Xx)$, one has $\Divergm_\phi(\al|\be) \geq 0$. If $\phi$ is strictly convex, then one has $\Divergm_\phi(\al|\be)=0$ if and only if $\al=\be$.
\end{proposition}
\begin{proof}
Let $m=\al+\be$ and write $a=\d\al/\d m$ and $b=\d\be/\d m$. Then $a+b=1$ $m$-almost everywhere and, using the perspective $\psi$ above, \(\Divergm_\phi(\al|\be)=\int \psi(a,b)\d m.\)
Since $m(\Xx)=2$, the measure $m/2$ is a probability. Jensen's inequality, the $1$-homogeneity of $\psi$, and $\phi(1)=0$ give
\(\frac12\Divergm_\phi(\al|\be) \geq \psi\pa{\frac12\int a\d m,\frac12\int b\d m} =\psi(1/2,1/2)=0.\)
If $\phi$ is strictly convex, the map $s\mapsto\psi(s,1-s)$ is strictly convex on its effective domain: distinct points on the line $u+v=1$ cannot lie on the same ray, which is the only degeneracy of a strict perspective. Equality in Jensen therefore forces $a=1/2$, hence $a=b$ and $\al=\be$. For arbitrary finite nonnegative measures, $\phi\geq0$ also implies $\phi'_\infty\geq0$, so every term in~\eqref{eq-phi-div} is nonnegative.
\end{proof}
\paragraph{Classical examples and topology.}
\paragraph{\texorpdfstring{Main families of $\phi$-divergences.}{Main families of phi-divergences.}}
\(\phi_\gamma(s)=\frac{s^\gamma-\gamma s+\gamma-1}{\gamma(\gamma-1)} \qquad(\gamma\neq0,1)\)
\(\Hellinger(\al,\be) \eqdef\Divergm_{\phi_H}(\al|\be)^{1/2} =\norm{\sqrt{\rho_\al}-\sqrt{\rho_\be}}_{L^2(\lambda)}.\)
\[
\JS(\alpha,\beta)^2
=
\frac12\KL\!\left(\alpha\middle|\frac{\alpha+\beta}{2}\right)
+
\frac12\KL\!\left(\beta\middle|\frac{\alpha+\beta}{2}\right),
\]
and satisfies $0\leq\JS(\al,\be)^2\leq\log 2$. Its exact entropy generator is

\[
\phi_{\JS}(s)
=\frac12\left[s\log s-(s+1)\log\pa{\frac{s+1}{2}}\right],
\qquad
\Divergm_{\phi_{\JS}}(\al|\be)=\JS(\al,\be)^2.
\]
\paragraph{Variational dual formula.}
\begin{proposition}[Variational representation of a $\phi$-divergence]\label{prop-phi-div-dual}
Assume that $\Xx$ is compact and let $\al,\be\in\Mm_+(\Xx)$. Define the Legendre transform of the entropy, already extended by $+\infty$ on negative arguments, by \(\phi^*(s) \eqdef \sup_{r\geq0}\{sr-\phi(r)\}.\)
Then
\eql{\label{eq-dual-div}
\Divergm_\phi(\al|\be)
=
\sup_{f\in\Cc(\Xx)}
\left\{\int_\Xx f\d\al-\int_\Xx\phi^*(f)\d\be\right\}.
}
Equivalently, the convex conjugate of $\Divergm_\phi(\cdot|\be)$ on $\Mm(\Xx)$ is
\eql{\label{eq-legendre}
\foralls f \in \Cc(\Xx), \quad
\Divergm_\phi^*(f|\be) = \int_\Xx \phi^*(f(x)) \d\be(x).
}
\end{proposition}
\begin{proof}
First suppose that $\phi'_\infty=+\infty$. Finite divergence then forces $\al=\rho\be$ for some $\rho\geq0$, and pointwise convex duality gives
\begin{align*}
\Divergm_\phi^*(f|\be)
&= \sup_{\rho\geq0}
\int_\Xx \bigl(f(x)\rho(x)-\phi(\rho(x))\bigr)\d\be(x)\\
&= \int_\Xx \sup_{r\geq0}\bigl(f(x)r-\phi(r)\bigr)\d\be(x)
= \int_\Xx \phi^*(f(x))\d\be(x).
\end{align*}
For finite $\phi'_\infty$, the effective domain of $\phi^*$ is bounded above by $\phi'_\infty$. Hence a singular mass contributes at most $\phi'_\infty\al^\perp(\Xx)$ to the pairing, exactly matching the recession term in~\eqref{eq-phi-div}; approximation by continuous test functions recovers this value. The same conjugate formula therefore holds in general. Proposition~\ref{prop-basic-phi-divergence-properties} gives convexity and weak-* lower semicontinuity, so Fenchel--Moreau yields~\eqref{eq-dual-div}.
\end{proof}
\subsection{GANs via Duality}
\label{sec-gan-duality}
\paragraph{Divergence-based adversarial losses.}
\eq{
\umin{\th} \Divergm_\phi(\al_\th|\be)
= \umin{\th} \usup{\f} \int_\Xx f(x) \d\al_\th(x) - \Divergm_\phi^*(f|\be)
= \umin{\th} \usup{\f} \int_\Xx f(g_\th(z)) \d\zeta(z) -
\frac{1}{m}\sum_{j=1}^m \phi^*(f(y_j)).
}
\(\min_\theta\max_\xi \int_\Zz f_\xi(g_\theta(z))\d\zeta(z) - \frac1m\sum_{j=1}^m \phi^*(f_\xi(y_j)).\)
For fixed $\theta$, this restriction gives a lower bound on the exact divergence unless the discriminator class is rich enough to attain the supremum. This distinction is essential for empirical data: if $\be$ is discrete and $\al_\theta$ is non-atomic, a superlinear divergence is $+\infty$, whereas its neural restricted objective can remain finite.

\(\widehat\phi_{\JS}(s)=s\log s-(s+1)\log\frac{s+1}{2}, \qquad \widehat\phi_{\JS}^*(u)=-\log(2-e^u),\quad u<\log2.\)
Thus $\Divergm_{\widehat\phi_{\JS}}=2\JS^2$, consistently with the normalization above. In practice the min--max problem is solved by alternating stochastic gradient descent/ascent on $(\theta,\xi)$.

\paragraph{Dual norms and integral probability metrics.}
\eq{
\umin{\th} \norm{\al_\th - \be}_B
= \umin{\th}
\usup{\f \in B} \int_{\X} \f(x) \d(\al_\th-\be)(x)
= \umin{\th}
\usup{\f \in B} \int_{\Zz} \f( g_\th(z)) \d\zeta - \frac{1}{m} \sum_{j=1}^m f(y_j).
}